%========================================================
% Experiment 1 — Single-Segment PMD
%========================================================

The first numerical experiment investigates the elementary mechanism by which
polarization-mode dispersion produces effective polarization decoherence for a
finite-bandwidth entangled-photon state. The objective is to isolate the
competition between spectral bandwidth and differential group delay (DGD) in a
configuration for which an analytical solution is available.

A single carrier-compensated birefringent segment is placed in arm \(A\), while
arm \(B\) is taken to be ideal. This provides a controlled connection between the
single-segment model of
Section~\ref{subsec:polarization_frequency_coupling} and the more general
segmented-fiber simulations considered later.

%--------------------------------------------------------
% Configuration and analytical prediction
%--------------------------------------------------------

\subsection{Configuration and Analytical Prediction}
\label{subsec:experiment_1_configuration_analytics}

The input polarization state is

\begin{equation}
|\Psi^+\rangle
=
\frac{1}{\sqrt{2}}
\left(
|01\rangle+|10\rangle
\right),
\qquad
\rho_{\mathrm{in}}
=
|\Psi^+\rangle\langle\Psi^+|.
\label{eq:experiment_1_input_state}
\end{equation}

The segment in arm \(A\) has birefringence axis

\[
\hat{\mathbf n}_A=\hat{\mathbf z},
\]

and its carrier-frequency differential phase is compensated,

\[
\Delta\phi_A(\omega_{0,A})=0.
\]

Introducing the frequency offset

\[
\Omega_A=\omega_A-\omega_{0,A},
\]

the residual transformation is therefore

\begin{equation}
U_A^{\mathrm{res}}
\left(
\omega_{0,A}+\Omega_A
\right)
=
\exp
\left[
\frac{i}{2}
\Omega_A
\Delta\tau_A
\sigma_z
\right],
\label{eq:experiment_1_single_segment_operator}
\end{equation}

with

\[
\tau_{\mathrm{DGD}}
=
|\Delta\tau_A|.
\]

Arm \(B\) undergoes the identity transformation. The zero carrier phase ensures
that the experiment isolates frequency-induced decoherence without introducing an
additional coherent rotation that could independently modify the fidelity.

The source is represented by the Gaussian joint spectral intensity

\begin{equation}
J(\Omega_A,\Omega_B)
\propto
\exp
\left[
-
\frac{(\Omega_A+\Omega_B)^2}{2\sigma_{\Sigma}^2}
-
\frac{(\Omega_A-\Omega_B)^2}{2\sigma_{\Delta}^2}
\right].
\label{eq:experiment_1_gaussian_jsi}
\end{equation}

For this experiment,

\begin{equation}
\sigma_{\Sigma}
=
\sigma_{\Delta}
=
\sqrt{2}\sigma_\omega,
\label{eq:experiment_1_uncorrelated_spectrum}
\end{equation}

so that the JSI factorizes and both photons have independent Gaussian marginal
spectra with root-mean-square width \(\sigma_\omega\).

For photon \(A\), the spectral trace produces the coherence factor derived in
Section~\ref{subsec:polarization_frequency_coupling},

\begin{equation}
\kappa
=
\exp
\left[
-\frac{1}{2}
\sigma_\omega^2
\tau_{\mathrm{DGD}}^2
\right].
\label{eq:experiment_1_coherence_factor}
\end{equation}

Introducing the dimensionless PMD strength

\begin{equation}
\xi
=
\sigma_\omega\tau_{\mathrm{DGD}},
\label{eq:experiment_1_dimensionless_pmd_strength}
\end{equation}

the result becomes

\begin{equation}
\boxed{
\kappa(\xi)=e^{-\xi^2/2}
}.
\label{eq:experiment_1_kappa_xi}
\end{equation}

The reduced polarization state is consequently

\begin{equation}
\rho_{\mathrm{pol}}^{\mathrm{out}}
=
\frac{1}{2}
\begin{pmatrix}
0 & 0 & 0 & 0 \\
0 & 1 & \kappa & 0 \\
0 & \kappa & 1 & 0 \\
0 & 0 & 0 & 0
\end{pmatrix},
\label{eq:experiment_1_analytical_density_matrix}
\end{equation}

with correlation tensor

\begin{equation}
T
=
\begin{pmatrix}
\kappa & 0 & 0 \\
0 & \kappa & 0 \\
0 & 0 & -1
\end{pmatrix}.
\label{eq:experiment_1_analytical_tensor}
\end{equation}

The expected state metrics are therefore

\begin{equation}
\gamma
=
\frac{1+\kappa^2}{2},
\qquad
F_{\Psi^+}
=
\frac{1+\kappa}{2},
\qquad
C=\kappa,
\qquad
\tau=\kappa^2,
\label{eq:experiment_1_analytical_metrics}
\end{equation}

and

\begin{equation}
S_{\max}
=
2\sqrt{1+\kappa^2}.
\label{eq:experiment_1_analytical_chsh}
\end{equation}

Thus, for \(\xi=0\), the original pure maximally entangled state is recovered,
whereas for \(\xi\rightarrow\infty\),

\[
C\rightarrow0,
\qquad
\gamma\rightarrow\frac{1}{2},
\qquad
F_{\Psi^+}\rightarrow\frac{1}{2},
\qquad
S_{\max}\rightarrow2.
\]

%--------------------------------------------------------
% Numerical procedure
%--------------------------------------------------------

\subsection{Numerical Procedure}
\label{subsec:experiment_1_numerical_procedure}

The numerical implementation evaluates the frequency-dependent propagation on a
discrete grid of angular-frequency offsets. For every value of
\(\sigma_\omega\), the interval

\[
\Omega
\in
[-6\sigma_\omega,+6\sigma_\omega]
\]

is sampled using \(81\) uniformly spaced nodes per arm. The corresponding
normalized joint spectral weights \(Q_{kl}\) are constructed from
Eq.~\eqref{eq:experiment_1_gaussian_jsi}.

For every spectral pair, the propagated polarization state is

\begin{equation}
\rho_{kl}
=
\left[
U_A(\omega_{0,A}+\Omega_{A,k})
\otimes I_2
\right]
\rho_{\mathrm{in}}
\left[
U_A(\omega_{0,A}+\Omega_{A,k})
\otimes I_2
\right]^\dagger,
\end{equation}

and the unresolved spectrum is traced out numerically through

\begin{equation}
\rho_{\mathrm{pol}}^{\mathrm{num}}
=
\sum_{k,l}
Q_{kl}\rho_{kl}.
\label{eq:experiment_1_numerical_spectral_trace}
\end{equation}

The parameter range used in the experiment is summarized in
Table~\ref{tab:experiment_1_parameters}.

\begin{table}[htbp]
    \centering
    \caption{Numerical parameters used in Experiment~1.}
    \label{tab:experiment_1_parameters}

    \begin{tabular}{lll}
        \hline
        Parameter & Range or value & Sampling \\
        \hline
        DGD
        & \(0\)--\(20~\mathrm{ps}\)
        & \(61\) points \\
        \(\sigma_\omega\)
        & \(2.5\times10^{10}\)--\(2.0\times10^{11}~\mathrm{rad/s}\)
        & \(21\) points \\
        \(\sigma_f=\sigma_\omega/(2\pi)\)
        & approximately \(4\)--\(32~\mathrm{GHz}\)
        & derived \\
        Spectral interval
        & \(\pm6\sigma_\omega\)
        & \(81\) nodes per arm \\
        \(\Delta\phi_A(\omega_{0,A})\)
        & \(0\)
        & fixed \\
        \(\hat{\mathbf n}_A\)
        & \(\hat{\mathbf z}\)
        & fixed \\
        \hline
    \end{tabular}
\end{table}

The resulting sweep contains \(1281\) bandwidth--DGD combinations and reaches

\[
\xi_{\max}
=
(2.0\times10^{11})
(20\times10^{-12})
=
4,
\]

thereby covering weak, intermediate, and strong polarization-coherence
suppression.

%--------------------------------------------------------
% Results
%--------------------------------------------------------

\subsection{Results and Analytical Comparison}
\label{subsec:experiment_1_results}

Figure~\ref{fig:experiment_1_single_segment_pmd_regime_map} shows the numerical
dependence of concurrence, global purity, Bell-state fidelity, and maximal CHSH
parameter on DGD and spectral bandwidth. The bandwidth is expressed as

\[
\sigma_f
=
\frac{\sigma_\omega}{2\pi}
\]

in gigahertz.

\begin{figure}[htbp]
    \centering
    \includegraphics[
        width=\textwidth
    ]{
        figures/experiment_1/experiment_1_single_segment_pmd_regime_map.png
    }

    \caption{
        Bandwidth--DGD regime maps for a single carrier-compensated
        birefringent segment acting on arm \(A\). The dashed curve denotes
        \(\xi=\sigma_\omega\tau_{\mathrm{DGD}}=1\).
        Increasing either spectral bandwidth or DGD produces the same
        progressive loss of unresolved polarization coherence through their
        dimensionless product \(\xi\).
    }
    \label{fig:experiment_1_single_segment_pmd_regime_map}
\end{figure}

The numerical results reproduce the expected dependence on
\(\xi=\sigma_\omega\tau_{\mathrm{DGD}}\). In the weak-PMD region
\(\xi\ll1\), the polarization transformation varies negligibly over the occupied
spectrum and the state remains approximately pure and maximally entangled. As
\(\xi\) increases, different spectral components acquire increasingly different
relative phases. After the spectral trace, the transverse coherence decreases,
leading simultaneously to reduced purity, fidelity, and concurrence.

The dashed contour \(\xi=1\) identifies the region in which the DGD becomes
comparable with the inverse spectral bandwidth. It is not a sharp threshold, but
provides a useful reference between weak and substantial PMD-induced coherence
reduction.

In the strong-PMD limit, the state approaches

\begin{equation}
\rho_{\mathrm{pol}}^{\mathrm{out}}
\longrightarrow
\frac{1}{2}
\left(
|01\rangle\langle01|
+
|10\rangle\langle10|
\right),
\end{equation}

so that the transverse correlations vanish while

\[
T_{zz}=-1
\]

remains unchanged. The state therefore retains classical anticorrelation in the
computational basis even after the polarization entanglement has disappeared.

For this particular state family,

\[
S_{\max}
=
2\sqrt{1+C^2},
\]

and therefore \(S_{\max}>2\) for every finite value of \(\xi\), approaching the
classical limit only asymptotically as \(C\rightarrow0\). This behavior is a
property of the ideal single-axis dephasing model and should not be generalized
to more complex fiber channels.

The point-by-point numerical values were also compared with the analytical
predictions. The maximum absolute discrepancies over the complete sweep are
reported in Table~\ref{tab:experiment_1_numerical_errors}.

\begin{table}[htbp]
    \centering
    \caption{
        Maximum absolute numerical--analytical discrepancies in
        Experiment~1.
    }
    \label{tab:experiment_1_numerical_errors}

    \begin{tabular}{lc}
        \hline
        Quantity & Maximum absolute error \\
        \hline
        Purity & \(3.512\times10^{-9}\) \\
        Fidelity & \(1.978\times10^{-9}\) \\
        Concurrence & \(3.956\times10^{-9}\) \\
        Tangle & \(7.025\times10^{-9}\) \\
        Entanglement of formation & \(5.471\times10^{-9}\) \\
        \(S_{\max}\) & \(5.246\times10^{-9}\) \\
        \(T_{xx}\) & \(3.956\times10^{-9}\) \\
        \(T_{yy}\) & \(3.956\times10^{-9}\) \\
        \(T_{zz}\) & \(1.332\times10^{-15}\) \\
        \hline
    \end{tabular}
\end{table}

The agreement at approximately \(10^{-9}\) or better confirms the numerical
implementation of the frequency-dependent segment operator, spectral
discretization, spectral partial trace, and subsequent state characterization.
The remaining discrepancy is consistent with the finite spectral interval and
frequency-grid discretization.

%--------------------------------------------------------
% Physical interpretation
%--------------------------------------------------------

\subsection{Physical Interpretation}
\label{subsec:experiment_1_physical_interpretation}

The experiment shows that DGD alone is not sufficient to determine the influence
of PMD on an entangled photon. The same differential delay may produce negligible
degradation for a narrow spectrum and substantial degradation for a broader one.
The relevant scale is instead

\[
\boxed{
\xi=\sigma_\omega\tau_{\mathrm{DGD}}
}.
\]

This result also establishes a direct connection with the effective dephasing model
introduced in Section~\ref{sec:statistical_phase_fluctuations}. In the present
physical PMD model, the effective coherence parameter is not assigned
phenomenologically but follows from measurable source and propagation parameters,

\begin{equation}
|\mu|
\longleftrightarrow
|\kappa|
=
\exp
\left[
-\frac{1}{2}
\sigma_\omega^2
\tau_{\mathrm{DGD}}^2
\right].
\label{eq:experiment_1_effective_dephasing_connection}
\end{equation}

Experiment~1 therefore provides the first quantitative connection between the
polarization-only effective models and frequency-dependent fiber propagation.
Its assumptions are intentionally restrictive: PMD acts in only one arm, the
birefringence axis is fixed, and the photon frequencies are uncorrelated. These
restrictions are relaxed in the following experiment, where PMD acts in both arms
and the role of the joint spectral correlations is investigated.